\SetPageTheme{story}
\PageHeader{Lectia 1}{Ce fac operatorii \texttt{/} si \texttt{\%}}{L1 • P1}

\begin{missionbox}[Povestea operatorilor]
  Daca vrem sa protejam sau sa verificam o informatie, calculatorul nu poate lucra doar cu impresii generale. El trebuie sa faca operatii precise.

  Uneori are nevoie sa afle ce parte intreaga obtine dupa o impartire. Alteori are nevoie exact de rest. In problemele de securitate, in codurile de verificare sau in descompunerea numerelor, restul nu este deloc un detaliu: el poate fi chiar informatia-cheie.

  De aceea, inainte sa invatam algoritmul principal, trebuie sa intelegem foarte bine aceste doua operatii.
\end{missionbox}

\begin{conceptbox}[Doua rezultate diferite]
  Pentru numere intregi:
  \begin{itemize}
    \item \texttt{a / b} inseamna catul impartirii;
    \item \texttt{a \% b} inseamna restul impartirii.
  \end{itemize}
\end{conceptbox}

\begin{guidebox}[Exercitiu ghidat]
  Completeaza:
  \begin{itemize}
    \item \texttt{17 / 5 =} \hrulefill
    \item \texttt{17 \% 5 =} \hrulefill
  \end{itemize}
\end{guidebox}

\newpage
\SetPageTheme{guided}
\PageHeader{Lectia 1}{Exersam operatorii pas cu pas}{L1 • P2}

\begin{examplebox}[Exemple explicate]
  \begin{itemize}
    \item \texttt{25 / 4 = 6} pentru ca 4 intra de 6 ori in 25;
    \item \texttt{25 \% 4 = 1} pentru ca ramane 1;
    \item \texttt{48 / 10 = 4};
    \item \texttt{48 \% 10 = 8}.
  \end{itemize}
\end{examplebox}

\begin{guidebox}[Set ghidat]
  Completeaza valorile:
  \begin{itemize}
    \item \texttt{31 / 6 =} \hrulefill
    \item \texttt{31 \% 6 =} \hrulefill
    \item \texttt{92 / 10 =} \hrulefill
    \item \texttt{92 \% 10 =} \hrulefill
    \item \texttt{104 / 10 =} \hrulefill
    \item \texttt{104 \% 10 =} \hrulefill
  \end{itemize}
\end{guidebox}

\begin{guidebox}[Intrebare de intelegere]
  De ce pot doua operatii diferite aplicate aceluiasi numar sa dea doua informatii utile in acelasi timp?
  \AnswerSpace{3}
\end{guidebox}

\newpage
\SetPageTheme{guided}
\PageHeader{Lectia 1}{Trecem spre autonomie}{L1 • P3}

\begin{solobox}[Exercitii semi-ghidate]
  1. Completeaza:
  \begin{itemize}
    \item \texttt{63 / 8 =} \hrulefill
    \item \texttt{63 \% 8 =} \hrulefill
  \end{itemize}

  2. Scrie cuvinte potrivite:
  \begin{itemize}
    \item operatorul \texttt{/} ne spune \hrulefill
    \item operatorul \texttt{\%} ne spune \hrulefill
  \end{itemize}
\end{solobox}

\begin{guidebox}[Observatie]
  Daca restul este 0, spunem ca primul numar este divizibil cu al doilea.
\end{guidebox}

\begin{solobox}[Exercitiu independent scurt]
  Alege trei exemple proprii de impartiri si scrie pentru fiecare atat catul, cat si restul.
  \AnswerSpace{4}
\end{solobox}

\newpage
\SetPageTheme{challenge}
\PageHeader{Lectia 1}{Lucru independent}{L1 • P4}

\begin{solobox}[Fisa independenta]
  1. Calculeaza:
  \begin{itemize}
    \item \texttt{77 / 9}
    \item \texttt{77 \% 9}
    \item \texttt{145 / 12}
    \item \texttt{145 \% 12}
  \end{itemize}

  2. Explica de ce restul este intotdeauna mai mic decat impartitorul.

  3. Da un exemplu de situatie in care nu te intereseaza catul, ci numai restul.
  \AnswerSpace{8}
\end{solobox}
